% !TeX root = ../../Book1.tex
\section{诱导内积、Gram 行列式与体积}
\label{b1:sec:B103}

到目前为止，外积只记录代数上的方向和变号规则。它没有长度，因而也没有面积。现在给 \(V\) 选定欧氏内积 \(\langle\cdot,\cdot\rangle_V\)，并记向量长度为 \(|v|\)。本节把这个内积传到各次外空间；传递规则必须使一组标准正交向量张成的单位平行多面体仍有单位体积，而不能因外积归一化而多出一个阶乘。

\subsection{外空间上的欧氏内积}
\label{b1:subsec:B10301}

\begin{theorem}[外空间的诱导内积]\label{b1:thm:app-b-gram-inner-product}
存在唯一的内积，使任意标准正交基 \(e_1,\ldots,e_n\) 所给出的 \(e_I\) 都成为 \(\bigwedge^kV\) 的标准正交基。对任意 \(v_1,\ldots,v_k,w_1,\ldots,w_k\in V\)，有
\begin{equation}\label{b1:eq:app-b-exterior-inner-product}
 \left\langle
    v_1\wedge\cdots\wedge v_k,\,
    w_1\wedge\cdots\wedge w_k
 \right\rangle
  =\det\bigl(\langle v_i,w_j\rangle_V\bigr)_{i,j=1}^k.
\end{equation}
特别地，
\[
 |v_1\wedge\cdots\wedge v_k|
  =\sqrt{\det\bigl(\langle v_i,v_j\rangle_V\bigr)_{i,j=1}^k}.
\]
零次空间 \(\bigwedge^0V=\mathbb R\) 取通常的实内积。
\end{theorem}
\begin{proof}
先选定一组标准正交基，声明 \(e_I\) 标准正交，就得到一个正定内积。若 \(A,B\) 分别是 \(v_i,w_j\) 的列坐标矩阵，则外积的子式展开给出
\[
 \left\langle\sum_I\det(A_I)e_I,\sum_J\det(B_J)e_J\right\rangle
   =\sum_I\det(A_I)\det(B_I)
   =\det(A^{\mathsf T}B).
\]
最后一个等式是\subsecref{b1:subsec:170104}证明的 Cauchy--Binet 公式；右侧矩阵的第 \(i,j\) 项正是 \(\langle v_i,w_j\rangle_V\)，从而得到所述公式。

公式右侧与最初所选的标准正交基无关，而所有可分解外向量张成 \(\bigwedge^kV\)。所以该公式唯一决定整个双线性型，也证明换用另一组标准正交基会得到同一内积。将 \(w_i=v_i\) 并取非负平方根，便得到范数公式。零次空间的空行列式取 \(1\)，与约定一致。
\end{proof}

Simon 的《Lectures on Geometric Measure Theory》\cite[\S25, pp. 129--131]{Simon1983Lectures}在介绍流之前采用了这一标准正交基约定；本节的证明补出由外积坐标到内积不变性的代数步骤。特别是，这里的 \(|\xi|\) 是外空间作为有限维欧氏空间的长度，不预先赋予“一般外向量所携带的总面积”这一含义。

在上一节给出的张量交错部分识别下，情况略有不同。如果张量空间中各个 \(e_{i_1}\otimes\cdots\otimes e_{i_k}\) 标准正交，那么
\[
 \left|
   \frac1{k!}\sum_{\sigma\in S_k}
       \operatorname{sgn}(\sigma)
       e_{\sigma(1)}\otimes\cdots\otimes e_{\sigma(k)}
 \right|^2=\frac1{k!}.
\]
它的外代数对应物 \(e_1\wedge\cdots\wedge e_k\) 的长度却是 \(1\)。两种空间之间的代数同构不自动是等距同构。若从张量子空间出发定义外空间内积，就必须作相应的归一化。

\begin{proposition}\label{b1:prop:app-b-decomposable-plane}
对 \(1\leq k\leq n\)，外积 \(v_1\wedge\cdots\wedge v_k\) 非零当且仅当各 \(v_i\) 线性无关。若
\(\xi=v_1\wedge\cdots\wedge v_k\ne0\)，则它唯一确定平面
\[
 P_\xi=\operatorname{span}\{v_1,\ldots,v_k\}
      =\{\,v\in V:v\wedge\xi=0\,\}.
\]
存在该平面的一组标准正交基 \(q_1,\ldots,q_k\)，使
\[
 \xi=|\xi|\,q_1\wedge\cdots\wedge q_k.
\]
\end{proposition}
\begin{proof}
若各向量线性相关，交错性使外积为零；若线性无关，将它们扩充成 \(V\) 的一组基，则其外积是该组基对应的一个外空间基向量，所以非零。

把 \(v_1,\ldots,v_k\) 扩充为基。属于其张成空间的向量与 \(\xi\) 外积为零；不属于这一空间的向量至少有一个补空间方向的非零系数，与 \(\xi\) 外积后得到互不相同的基外向量的非平凡线性组合，因而非零。这证明了平面的内在刻画。

在这个平面内选标准正交基 \(q_i\)，将 \(v_j=\sum_i a_{ij}q_i\) 代入外积，得到
\[
 \xi=(\det A)\,q_1\wedge\cdots\wedge q_k.
\]
由诱导范数，\(|\xi|=|\det A|\)。若行列式为负，就把 \(q_1\) 改为 \(-q_1\)，便得到正的系数。这个符号选择将在定向的定义中发挥作用。
\end{proof}

长度相同并不意味着代表同一个平面；同一个平面也可以有相反方向。例如 \(e_1\wedge e_2\) 与 \(-e_1\wedge e_2\) 张成同一二维平面、具有相同长度，却是不同外向量。外积因此同时保存了一个平面、一个尺度和一个方向，取绝对长度时才丢弃后两者中的方向信息。

\subsection{体积与非正交坐标}
\label{b1:subsec:B10302}

设 \(v_1,\ldots,v_k\) 线性无关，考虑其生成的平行多面体
\[
 Q=\left\{\sum_{j=1}^k t_jv_j:0\leq t_j\leq1\right\}.
\]
在平面 \(P=\operatorname{span}\{v_j\}\) 上使用由环境内积诱导的距离，并采用\chapref{b1:ch:15}的 Hausdorff 测度归一化。由线性面积公式，
\begin{equation}\label{b1:eq:app-b-parallelepiped-volume}
 \mathcal H^k(Q)
 =|v_1\wedge\cdots\wedge v_k|
 =\sqrt{\det\bigl(\langle v_i,v_j\rangle_V\bigr)}.
\end{equation}
若向量相关，则 \(Q\) 落在低维平面中，两侧都为零。这里先有外代数内积，再借\chapref{b1:ch:17}已证的测度公式解释其长度，因而没有用直观体积来反过来证明代数公式。

还可直接看到“底乘高”的结构。将 \(v_j\) 依次分解为其在前面向量张成空间中的投影与正交余量 \(w_j\)。在外积中，投影项含有重复方向而消失，所以
\[
 v_1\wedge\cdots\wedge v_k=w_1\wedge\cdots\wedge w_k,
 \quad
 |v_1\wedge\cdots\wedge v_k|=\prod_{j=1}^k|w_j|.
\]
每个 \(|w_j|\) 都是新边相对于已有底面的垂直高度。由 \(|w_j|\leq|v_j|\) 得到
\[
 |v_1\wedge\cdots\wedge v_k|\leq\prod_{j=1}^k|v_j|.
\]
若每个 \(v_j\ne0\)，且它们线性无关，那么等号成立恰好要求每一步投影为零，也就是各边两两正交；线性相关时左侧为零，若右侧非零便不可能取等。

在非标准正交基 \(b_1,\ldots,b_n\) 下，记
\[
 G=(\langle b_i,b_j\rangle_V)_{i,j=1}^n.
\]
若 \(v_j\) 的列坐标组成 \(n\times k\) 矩阵 \(A\)，则
\[
 |v_1\wedge\cdots\wedge v_k|^2=\det(A^{\mathsf T}GA).
\]
坐标中的 \(A^{\mathsf T}A\) 只有在 \(G=I_n\) 时才是正确的 Gram 矩阵。尤其对全维平行多面体，
\[
 \mathcal H^n(Q)=\sqrt{\det G}\,|\det A|.
\]
这不是新的测度归一化，而是坐标基本身所围出的单位盒子已经具有体积 \(\sqrt{\det G}\)。

例如在 \(\mathbb R^2\) 中取 \(b_1=(1,0)\)、\(b_2=(1,1)\)。向量 \(b_2\) 的坐标为 \((0,1)\)，但真实长度为 \(\sqrt2\)；由于
\[
 G=\begin{pmatrix}1&1\\1&2\end{pmatrix},
\]
正确的坐标长度是 \(\sqrt{x^{\mathsf T}Gx}\)，而不是坐标数组的通常长度。另一方面 \(\det G=1\)，这组基围出的平行四边形面积恰为 \(1\)。长度失真与总体积失真是不同问题。

\subsection{外积的估计与一般外向量}
\label{b1:subsec:B10303}

\begin{proposition}\label{b1:prop:app-b-wedge-estimates}
若 \(\xi\in\bigwedge^kV\)、\(\eta\in\bigwedge^\ell V\)，且 \(k+\ell\leq n\)，则
\[
 |\xi\wedge\eta|
 \leq\sqrt{\binom{k+\ell}{k}}\,
        |\xi|\cdot|\eta|.
\]
若 \(\xi\) 可分解，则有更强的估计
\[
 |\xi\wedge\eta|\leq|\xi|\cdot|\eta|.
\]
当 \(k+\ell>n\) 时，外积为零。
\end{proposition}
\begin{proof}
选标准正交基，写 \(\xi=\sum_I\xi_Ie_I\)、\(\eta=\sum_J\eta_Je_J\)。对每个大小为 \(k+\ell\) 的指标组 \(K\)，外积的 \(e_K\) 系数是
\[
 \sum_{\substack{I\cup J=K\\I\cap J=\varnothing}}
       s_{I,J}\xi_I\eta_J,\quad s_{I,J}\in\{-1,1\}.
\]
和式有 \(\binom{k+\ell}{k}\) 项。对这个有限和用 Cauchy--Schwarz，再对 \(K\) 求和，得到
\[
 |\xi\wedge\eta|^2
 \leq\binom{k+\ell}{k}
      \sum_{I\cap J=\varnothing}\xi_I^2\eta_J^2
 \leq\binom{k+\ell}{k}
      \left(\sum_I\xi_I^2\right)
      \left(\sum_J\eta_J^2\right).
\]

若 \(\xi\ne0\) 可分解，由上一命题可选择标准正交基使
\(\xi=|\xi|e_1\wedge\cdots\wedge e_k\)。展开 \(\eta\) 后，凡含有前 \(k\) 个指标的项都被外积消去，其余项仍对应不同的标准正交基外向量。因此
\[
 |\xi\wedge\eta|^2
   =|\xi|^2
      \sum_{J\cap\{1,\ldots,k\}=\varnothing}\eta_J^2
   \leq|\xi|^2\cdot|\eta|^2.
\]
零元素情形直接成立；次数大于维数时则由外空间为零得到结论。
\end{proof}

可分解条件在第二个估计中有实际作用。例如在 \(\mathbb R^6\) 中取
\[
 \xi=e_1\wedge e_2+e_3\wedge e_4+e_5\wedge e_6.
\]
有 \(|\xi|=\sqrt3\)，而
\[
 \xi\wedge\xi
  =2(e_1\wedge e_2\wedge e_3\wedge e_4
     +e_1\wedge e_2\wedge e_5\wedge e_6
     +e_3\wedge e_4\wedge e_5\wedge e_6)
\]
的长度为 \(2\sqrt3>3=|\xi|^2\)。所以外空间的欧氏长度不能不分对象地当成一个使乘法常数总为 \(1\) 的范数。

同样，在 \(\mathbb R^4\) 中，\(\xi=e_1\wedge e_2+e_3\wedge e_4\) 的欧氏长度为 \(\sqrt2\)，却不能据此说“两张互相正交的单位面片总面积为 \(\sqrt2\)”。这个代数和只保留一个外向量；几何上的两块面片若作为不同的积分对象相加，还必须说明支集、重数以及怎样测试它们。\appref{b1:app:D}将据此引入流的质量和余质量；在可分解外向量上有关长度相容，但一般外向量上的质量范数与本节欧氏范数不得混用。
